\documentclass[12pt]{article}
\usepackage[utf8]{inputenc}
\usepackage[T1]{fontenc}
\usepackage{hyperref}
\title{Zoogocho Zapotec Language Resources for K-12 Curriculum Development}
\author{Compiled by Claude Sonnet 5.5 (claude-sonnet-5-5) for Andr\'as Kornai}
\date{}
\begin{document}
\maketitle

\section{Introduction}
Zoogocho Zapotec (ISO 639-3 zpq; Glottocode zoog1238) is an Otomanguean language of Oaxaca, Mexico, with about 1,000 speakers. It is spoken in San Bartolom\'e Zoogocho, Santa Mar\'ia Yalina, Tabehua and Oaxaca City. The main materials found are a bilingual dictionary and a reference grammar.

\section{Primary Educational Resources}
No primers, textbooks or school programmes were found in the search results.

\section{Linguistic and Reference Materials}
\begin{itemize}
\item Rebecca Long C. and Sofronio Cruz M., \emph{Diccionario zapoteco de San Bartolom\'e Zoogocho, Oaxaca}, a Zapotec--Spanish dictionary of about 530 pages (library record): \url{https://si.edu/object/siris_sil_568363}
\item Aron Huey Sonnenschein, \emph{A Descriptive Grammar of San Bartolom\'e Zoogocho Zapotec} (Lincom, 2005), bibliographic record: \url{https://glottolog.org/resource/reference/id/124784}
\item Sonnenschein's 2004 USC dissertation of the same title, in the USC Digital Library: \url{https://digitallibrary.usc.edu/asset-management/2A3BF1LQ0DL2}
\end{itemize}

\section{Religious and Community Resources}
None were found in the search results.

\section{Digital and Audio Resources}
No audio, video or corpus was found. The dissertation reportedly analyses three Zoogocho texts, which may be a source of text material.

\section{Curriculum Development Context}
The dictionary and grammar are useful bases for lexicon and orthography work. No classroom materials were found, and the language is not known to be taught in schools. Curriculum development would depend on community partnership.
\end{document}